\chapter{Complementary Code Listings}
\label{chap:complementary-codes}

This appendix provides representative code excerpts that complement the tool and benchmark
descriptions in the main chapters. They illustrate the formats consumed and produced by our
tools and the firmware structure of the benchmark scenarios.

\lstset{style=mystyle}

% ====================================================================
\section{IAR EWARM linker map excerpt}
\label{sec:appendix-linker-map}

Listing~\ref{lst:linker-map} shows an abbreviated IAR EWARM linker map for a USB HID
project. This is the format parsed by MCU Workflow Studio
(Chapter~\ref{chap:footprint-tool}). Each entry lists a module, its read-only code size,
read-only data, and read-write data --- the three quantities from which flash and RAM
footprint are derived.

\begin{lstlisting}[language={}, caption={Abbreviated IAR EWARM linker map (USB HID, STM32U575).}, label={lst:linker-map}]
*******************************************************************************
*** MODULE SUMMARY
***

    Module               ro code  ro data  rw data
    ------               -------  -------  -------
    main.o                   164       --       --
    stm32u5xx_hal.o          212       --        4
    stm32u5xx_hal_gpio.o     456       --       --
    stm32u5xx_hal_rcc.o    4 560       16       --
    stm32u5xx_hal_pcd.o    2 496       --       --
    stm32u5xx_ll_usb.o     2 492       --       --
    usbd_core.o            1 048       --       --
    usbd_ctlreq.o            876       --       --
    usbd_hid.o               620       --       --
    usbd_conf.o              524       --       32
    usbd_desc.o              461      268       --
    system_stm32u5xx.o       128       16        4
    startup_stm32u575xx.o    394      512       --
    ------               -------  -------  -------
    Total:               14 431      812       40

*******************************************************************************
*** ENTRY LIST
***

    Entry                 Address    Size  Type  Object
    -----                 -------    ----  ----  ------
    HAL_Init             0x0800'1234   52  Code  stm32u5xx_hal.o
    HAL_RCC_OscConfig    0x0800'2000  1024  Code  stm32u5xx_hal_rcc.o
    HAL_RCC_ClockConfig  0x0800'2400   896  Code  stm32u5xx_hal_rcc.o
    USBD_Init            0x0800'3000   128  Code  usbd_core.o
    USBD_HID_Init        0x0800'3200    96  Code  usbd_hid.o
    ...
\end{lstlisting}

\pagebreak

% ====================================================================
\section{ST--TI functional footprint reproducibility record}
\label{sec:appendix-ti-functional-footprint}

This appendix records the build and accounting conditions behind the ST-versus-TI functional
footprint results in Chapter~\ref{chap:vs-ti}. The TI projects are MSPM33 SDK driverlib
examples~\cite{ti-mspm33-sdk}; their counterparts are STM32CubeC5 LL examples~\cite{stm32cubec5-sw}.
All twelve pairs were rebuilt with IAR EWARM~9.60.3 using high optimisation for size
(\texttt{CCOptLevel=3}, \texttt{CCOptStrategy=1}). The IAR linker map reports read-only code,
read-only data, and read-write data~\cite{iar-ewarm}. We use the following whole-image accounting:

\[
\text{Flash} = \text{RO code} + \text{RO data}, \qquad
\text{RAM} = \text{RW data}.
\]

Both projects in every pair reserve a 512-byte stack in the linker map. The table is therefore
reproducible from the matching IAR map files and includes the deployable image's startup and
vector-table cost. It is a portfolio of independent images, not the footprint of one firmware
that combines all scenarios.

\begin{table}[H]
\centering
\scriptsize
\caption{Reproducibility record for the ST--MSPM33 high/size functional footprint study (bytes).}
\label{tab:appendix-ti-functional-footprint}
\resizebox{\textwidth}{!}{%
\begin{tabular}{l r r r r r r}
\toprule
\textbf{Scenario} & \textbf{MSPM33 flash} & \textbf{STM32 flash} & \textbf{Delta} & \textbf{MSPM33 RAM} & \textbf{STM32 RAM} & \textbf{Delta} \\
\midrule
ADC timer trigger & 1\,476 & 1\,416 & -60 & 513 & 513 & 0 \\
GPIO toggle & 656 & 772 & +116 & 512 & 512 & 0 \\
I\textsuperscript{2}C DMA controller & 1\,928 & 1\,840 & -88 & 553 & 517 & -36 \\
I\textsuperscript{2}C DMA responder & 1\,744 & 1\,687 & -57 & 593 & 556 & -37 \\
I\textsuperscript{2}C interrupt controller & 1\,580 & 1\,732 & +152 & 557 & 514 & -43 \\
I\textsuperscript{2}C interrupt responder & 1\,388 & 1\,460 & +72 & 600 & 553 & -47 \\
I\textsuperscript{2}C polling controller & 1\,420 & 1\,376 & -44 & 556 & 513 & -43 \\
I\textsuperscript{2}C polling responder & 1\,412 & 1\,340 & -72 & 608 & 552 & -56 \\
SPI DMA controller & 1\,704 & 2\,151 & +447 & 592 & 556 & -36 \\
SPI DMA peripheral & 1\,876 & 1\,940 & +64 & 592 & 552 & -40 \\
UART printf & 1\,326 & 1\,318 & -8 & 512 & 512 & 0 \\
USART loopback & 1\,588 & 1\,580 & -8 & 532 & 532 & 0 \\
\midrule
\textbf{Total, 12 independent images} & \textbf{18\,098} & \textbf{18\,612} & \textbf{+514} & \textbf{6\,720} & \textbf{6\,382} & \textbf{-338} \\
\bottomrule
\end{tabular}%
}
\end{table}

Negative deltas favour STM32C5. The corresponding interpretation, fairness limits, and
scenario-flow diagrams are presented in Section~\ref{sec:ti-functional-analysis}; this appendix
is intentionally limited to reproducibility facts and raw whole-image totals.

\pagebreak

% ====================================================================
\section{ST--TI FreeRTOS reproducibility record}
\label{sec:appendix-ti-freertos}

The FreeRTOS middleware benchmark in Section~\ref{sec:ti-freertos} uses three matched project
pairs from STM32CubeC5 and MSPM33 SDK~1.04.00.00~\cite{stm32cubec5-sw,ti-mspm33-sdk}. All six
images were built with IAR EWARM~9.60.3 at high optimisation for size
(\texttt{CCOptLevel=3}, \texttt{CCOptStrategy=1})~\cite{iar-ewarm}. The common configuration
includes a 1\,kHz tick, preemption enabled, time slicing disabled, dynamic allocation, ten
priorities, and an 8\,KiB FreeRTOS heap. STM32C5 was configured at 144\,MHz and MSPM33 at
32\,MHz; these are source configuration values rather than independent clock measurements.
Project manifests identify FreeRTOS Kernel~V11.2.0 for STM32C5 and V11.1.0 for MSPM33. In the
preemptive scenario, STM32 uses \texttt{vTaskSuspend()}/\texttt{vTaskResume()}, whereas MSPM33
uses direct task notifications and a final acknowledgement.

\begin{table}[H]
\centering
\caption{Raw UART observation windows for the ST--TI FreeRTOS benchmark.}
\label{tab:appendix-ti-freertos-runtime}
\small
\begin{tabularx}{\textwidth}{l l X}
\toprule
\textbf{Scenario} & \textbf{Platform} & \textbf{Recorded 30-second deltas} \\
\midrule
Basic & STM32C5 & 1\,295\,358; 1\,295\,330; 1\,295\,351. Elapsed: 30\,000; 30\,005; 30\,006\,ms. \\
Basic & MSPM33 & 287\,141; 287\,135. Elapsed: 30\,000; 30\,005\,ms. \\
Cooperative & STM32C5 & 33\,998\,170; 33\,997\,657 total worker iterations \\
Cooperative & MSPM33 & 5\,367\,603; 5\,367\,557; 5\,367\,559; 5\,367\,490 total worker iterations \\
Preemptive & STM32C5 & 8\,252\,379; 8\,252\,277 total completed chains \\
Preemptive & MSPM33 & 1\,066\,661; 1\,066\,639; 1\,066\,639; 1\,066\,637; 1\,066\,635; 1\,066\,655 total completed chains \\
\bottomrule
\end{tabularx}
\end{table}

\begin{table}[H]
\centering
\caption{IAR high/size map totals for the ST--TI FreeRTOS projects (bytes).}
\label{tab:appendix-ti-freertos-footprint}
\begin{tabular}{l r r r r}
\toprule
\textbf{Scenario} & \textbf{STM32 flash} & \textbf{MSPM33 flash} & \textbf{STM32 RAM} & \textbf{MSPM33 RAM} \\
\midrule
Basic & 9\,052 & 12\,854 & 10\,196 & 10\,140 \\
Cooperative & 9\,640 & 13\,430 & 9\,188 & 9\,132 \\
Preemptive & 10\,012 & 13\,894 & 9\,188 & 9\,112 \\
\midrule
\textbf{Total} & \textbf{28\,704} & \textbf{40\,178} & \textbf{28\,572} & \textbf{28\,384} \\
\bottomrule
\end{tabular}
\end{table}

The runtime record is not a controlled equal-clock experiment. Sample counts differ, no
warm-up policy is documented, and the preemptive implementations use different synchronisation
primitives. The tables preserve the raw evidence so that later equal-clock measurements can be
compared without replacing or silently reinterpreting the original observations.

\pagebreak

% ====================================================================
\section{Agent configuration file}
\label{sec:appendix-agent-config}

Listing~\ref{lst:agent-config} shows the configuration of the Creation-Porting specialist
agent from the MCU Benchmark Copilot Agent system
(Chapter~\ref{chap:benchmark-agent}). The YAML frontmatter declares metadata, description,
argument hints, and tool permissions; the Markdown body contains the behavioural
instructions, autonomy contract, and supported workflow directions.

\begin{lstlisting}[language={}, caption={Creation-Porting agent configuration (abbreviated).}, label={lst:agent-config}]
---
name: Creation-Porting
description: "Create benchmark projects from scratch or port existing
  benchmarks in any direction (STM32->competitor, competitor->STM32,
  vendor->vendor, board->board). Use when: generating target project
  files, creating from a scenario description, separating
  portable/RTOS/board logic, or acquiring missing reference firmware."
argument-hint: A benchmark requirement, scenario description, scenario
  name, firmware name, or a reference project plus target platform.
tools: ['read', 'search', 'edit', 'vscode', 'todo', 'web', 'execute']
---

# Creation-Porting Agent

## Supported Directions
- From scratch (scenario description only)
- STM32 -> competitor
- Competitor -> STM32
- Vendor A -> Vendor B
- Board A -> Board B
- Name only (resolve, acquire, create)

## Autonomy Contract
- Operates autonomously for file generation, project creation, and
  material acquisition (web fetches from official sources).
- Asks only before flashing/programming/resetting hardware.
- Acquires missing firmware, CMSIS, BSP, startup, linker from
  official vendor GitHub repos and SDKs without confirmation.

## IAR Template Resolution
Before creating any project:
1. Look up target board in config/iar_templates.yaml
2. Follow resolution: exact board -> family fallback -> type fallback
3. Acquire template if status == "needs_download"
4. Use materialised template as read-only project baseline
5. Place benchmark code through app_integration_points only

## Multi-Layer Separation
1. Portable application layer (writable) - benchmark logic
2. RTOS / middleware layer (protected) - select, don't modify
3. Board-specific layer (protected) - select existing vendor files
4. Build surface (configure only) - .ewp references, defines, paths

## Semantic Preservation Rule
Preserve semantics and observable behaviors, not vendor API names.
Map CMSIS-RTOS2 <-> native FreeRTOS <-> ThreadX as needed via
app-level adapter files. Never modify vendor HAL source.
\end{lstlisting}

\pagebreak

% ====================================================================
\section{FreeRTOS cooperative benchmark scenario}
\label{sec:appendix-freertos-coop}

Listing~\ref{lst:freertos-coop} shows the core task logic of the cooperative processing
benchmark scenario from Chapter~\ref{chap:vs-gd32}. Five equal-priority threads each
increment a counter and yield; a reporter thread prints the results every 30~seconds.

\begin{lstlisting}[language=C, caption={FreeRTOS cooperative scenario --- task functions (STM32, simplified).}, label={lst:freertos-coop}]
/* Shared counters for the five worker threads */
volatile uint32_t counters[5] = {0};

/* Worker task: increment own counter, then yield */
void WorkerTask(void *param)
{
    uint32_t id = (uint32_t)param;
    for (;;)
    {
        counters[id]++;
        taskYIELD();
    }
}

/* Reporter task: print all counters every 30 s */
void ReporterTask(void *param)
{
    (void)param;
    for (;;)
    {
        vTaskDelay(pdMS_TO_TICKS(30000));
        printf("Counters: %lu %lu %lu %lu %lu\r\n",
               counters[0], counters[1], counters[2],
               counters[3], counters[4]);
        /* Reset for next interval */
        for (int i = 0; i < 5; i++)
            counters[i] = 0;
    }
}

/* Task creation (in main, after HAL and kernel init) */
void CreateBenchmarkTasks(void)
{
    for (uint32_t i = 0; i < 5; i++)
    {
        xTaskCreate(WorkerTask, "Worker", 128,
                    (void *)i, tskIDLE_PRIORITY + 1, NULL);
    }
    xTaskCreate(ReporterTask, "Reporter", 256,
                NULL, tskIDLE_PRIORITY + 2, NULL);
    vTaskStartScheduler();
}
\end{lstlisting}

\pagebreak

% ====================================================================
\section{Footprint tool CLI output}
\label{sec:appendix-tool-output}

Listing~\ref{lst:tool-output} shows a representative terminal output from a footprint
benchmark run using the MCU Workflow Studio command-line interface
(Chapter~\ref{chap:footprint-tool}).

\begin{lstlisting}[language={}, caption={MCU Workflow Studio CLI --- benchmark summary output.}, label={lst:tool-output}]
$ mcu-workflow-cli benchmark \
    --source stm32_usb_hid.map --source-vendor stm32 \
    --target gd32_usb_hid.map  --target-vendor gd32 \
    --out results/

[INFO] Parsing source map: stm32_usb_hid.map (IAR EWARM)
[INFO] Parsing target map: gd32_usb_hid.map (IAR EWARM)
[INFO] Classifying 47 source symbols into layers...
[INFO] Classifying 31 target symbols into layers...
[INFO] Mapping engine: 42 pairs scored, 38 accepted (confidence >= 0.72)
[INFO] 4 pairs queued for review (confidence 0.55-0.72)

=== FOOTPRINT COMPARISON (USB HID) ===

Layer           | STM32 ROM | GD32 ROM  | Diff
----------------|-----------|-----------|--------
Application     |   2 580 B |   1 942 B | +33%
HAL drivers     |  10 426 B |     820 B | +1171%
MW stack        |   2 859 B |   5 743 B | -50%
----------------|-----------|-----------|--------
TOTAL ROM       |  16 269 B |   9 175 B | +77%
TOTAL RAM       |   3 446 B |   9 608 B | -65%

Verdict: STM32 uses +77% more ROM but -65% less RAM.
Largest classified STM32 modules:
    HAL RCC:       4560 B
    HAL/LL USB:    4988 B
Causal attribution: unresolved without complete maps and classification diagnostics.

[INFO] Results written to results/benchmark_results.xlsx
[INFO] Mapping diagnostics: results/mapping_diagnostics.json
[INFO] Review queue (4 items): results/review_queue.csv
\end{lstlisting}

\pagebreak

% ====================================================================
\section{Static-analysis automation script (ST vs.\ TI)}
\label{sec:appendix-analyze-firmware}

Listing~\ref{lst:analyze-firmware} shows an abbreviated version of the Python script used
to extract API surface and documentation coverage metrics for the ST-versus-TI static
analysis (Chapter~\ref{chap:vs-ti}). The full script analyses duplication, documentation,
and API design in a single pass over the driver source trees.

\begin{lstlisting}[language=Python, caption={Abbreviated \texttt{analyze\_firmware.py} --- API and documentation extraction.}, label={lst:analyze-firmware}]
"""
Firmware benchmark analysis script:
- Code duplication detection (sliding-window hash)
- Documentation coverage (Doxygen tag extraction)
- API surface area analysis (functions, macros, enums)
"""
import os, re, json, hashlib
from collections import defaultdict, Counter
from pathlib import Path

TI_DIR = r"mspm33_sdk/source/ti/driverlib"
ST_DIR = r"STM32CubeC5/stm32c5xx_drivers/ll"

def get_c_files(directory):
    files = []
    for root, _, filenames in os.walk(directory):
        for f in filenames:
            if f.endswith(('.c', '.h')):
                files.append(os.path.join(root, f))
    return files

def analyze_api(files):
    """Extract API surface metrics from C source files."""
    result = {
        'public_functions': 0, 'static_functions': 0,
        'inline_functions': 0, 'macros_functional': 0,
        'macros_constant': 0, 'enums': 0, 'typedefs': 0,
        'param_distribution': Counter(),
    }
    for filepath in files:
        with open(filepath, 'r', errors='ignore') as f:
            content = f.read()
        # Count public functions (not static)
        pub = re.findall(
            r'(?<!static\s)(?:__STATIC_INLINE\s+|inline\s+)?'
            r'(?:void|bool|uint\w+|int\w+|DL_\w+|LL_\w+)\s+'
            r'(\w+)\s*\(([^)]*)\)', content)
        for name, params in pub:
            result['public_functions'] += 1
            n = len([p for p in params.split(',')
                     if p.strip()]) if params.strip() else 0
            result['param_distribution'][n] += 1
        # Count enums, typedefs, macros
        result['enums'] += len(re.findall(
            r'\benum\s+\w*\s*\{', content))
        result['typedefs'] += len(re.findall(
            r'\btypedef\s+', content))
        result['macros_functional'] += len(re.findall(
            r'#define\s+\w+\([^)]*\)', content))
        result['macros_constant'] += len(re.findall(
            r'#define\s+\w+\s+[^(]', content))
    return result

if __name__ == '__main__':
    for label, path in [("TI", TI_DIR), ("ST", ST_DIR)]:
        files = get_c_files(path)
        api = analyze_api(files)
        print(f"\n=== {label} API Surface ===")
        for k, v in api.items():
            print(f"  {k}: {v}")
\end{lstlisting}
